\section{\ai{}}
\label{sec:benchmark}

\subsection{Task formulation}
\label{sec:formulation}

Each task presents an agent with a research repository, a base model to start from, and an inexpensive proxy metric that can be evaluated freely during development \S\ref{sec:tasks}. 
The agent is given 4 hours on one B300 GPU and a single objective: improve the training algorithm implemented in the repository.
Its submission is neither a number nor a trained model, but the repository’s source code after the agent’s modifications. 
After submission, the agent can no longer modify or interact with the code. The submitted repository is then run from scratch under a fixed budget, and the resulting model is scored by a predetermined evaluator (\S\ref{sec:protocol}).

% Notation. Introduced once, used for the rest of the paper. The point of writing
% it out is that "improvement" is a comparison between two runs of the same
% procedure differing only in the source code, and that is easier to see in
% symbols than in prose.
Formally, a task is a tuple $(C, a_0, q, m, d)$, where $C$ is the repository's source in its frozen
state, $a_0$ the model it starts from, $q$ an inexpensive proxy available to the agent, $m$
the final metric, and $d \in \{\uparrow, \downarrow\}$ its direction. An agent observes
$(C, a_0, q)$ under an exploration budget $T_{\mathrm{e}} = 4$ hours on one B300 GPU and returns a
rewritten source $C'$; it never evaluates $m$. Execution of $C'$ under a verification budget
$T_{\mathrm{v}} = 12$ hours yields a model $a(C')$, and the score of the resulting cell is
\begin{equation*}
  s(C') \;=\; m\bigl(a(C')\bigr),
\end{equation*}
where $m$ is computed by an evaluator $E$ that is fixed in advance and has no access to the agent's
execution environment. Applying the identical procedure to the unmodified source gives the baseline
$s(C) = m(a(C))$, and a submission constitutes an improvement precisely when
$s(C') \succ_d s(C)$, that is, $s(C') > s(C)$ when $d\,{=}\,\uparrow$ and $s(C') < s(C)$ when
$d\,{=}\,\downarrow$. The two executions differ in $C'$ against $C$ and in nothing else: the hardware, the budget, the evaluator and the evaluation asset are common to both.

What is held fixed is the measurement --- the evaluation asset, the final metric, and the evaluator
that computes it --- and what is open is the method. The agent may rewrite the training loop, the
objective, the optimizer, the data pipeline, the schedule, or all of them; the single line it may not
cross is the evaluation itself. Fixing only the outcome and its measurement asks the question this
paper is about: whether an agent can find a way to improve the system at all, and which part of the
system it reaches for when nothing constrains it.

% Terminology note (paragraph deleted on request): "system" = (model, harness,
% reasoning effort) and "cell" = one system on one task are still used in
% \S\ref{sec:protocol}, \S\ref{sec:families} and \S\ref{sec:results}. If they are
% never defined, define them at first use there instead.

\subsection{\ai{} algorithmic tasks}
\label{sec:tasks}

% The main task table for the rewritten Section 2. Starting models come from
% results/ai4ai_results.json (tasks[*].starting_artifact); metrics and directions
% from the final experiment sheet (build/lark/final_table.csv, rows 2-4). No
% reference column: baselines are defined in \S\ref{sec:baselines} and reported
% with the results, not here.
\begin{table}[t]
\centering
\caption{\textbf{The ten \ai{} tasks.} Each freezes a research repository and asks an agent to
improve the training algorithm that repository applies to its own model. \textbf{Starting model} is
the model $a_0$ the procedure begins with, and \textbf{Evaluation metric} the quantity a fixed
evaluator computes afterwards, with its direction. Between them the ten cover ten families of
training algorithm rather than ten instances of one, and their metrics are incommensurable.
$^{\dagger}$Weight averaging and one-shot pruning
do no training, and are executed once rather than trained to a horizon.}
\label{tab:tasks}
\small
\begin{adjustbox}{max width=\linewidth}
\begin{tabular}{llll}
\toprule
\textbf{Task} & \textbf{Algorithm family} & \textbf{Starting model} & \textbf{Evaluation metric} \\
\midrule
OpenR1     & supervised fine-tuning        & Qwen2.5-Coder-1.5B-Instruct   & LiveCodeBench $\uparrow$ \\
RAGEN      & multi-turn agentic RL         & Qwen2.5-3B-Instruct           & held-out solve rate $\uparrow$ \\
OPD        & on-policy distillation        & R1-Distill-Qwen-1.5B          & AIME 24/25 $\uparrow$ \\
BTRM       & Bradley--Terry reward model   & Mistral-7B-Instruct-v0.2      & RewardBench $\uparrow$ \\
DPO        & preference optimization       & merged Zephyr/Mistral-7B      & IFEval strict $\uparrow$ \\
DDPO       & diffusion RL                  & Stable Diffusion v1.5         & aesthetic score $\uparrow$ \\
NPO        & machine unlearning            & Llama-3.2-1B-Instruct         & balanced score $\uparrow$ \\
DiGress    & discrete graph diffusion      & QM9 graph diffusion model     & test NLL $\downarrow$ \\
\midrule
Model Soup$^{\dagger}$ & weight averaging  & 72 CLIP checkpoints           & ImageNet-V2 top-1 $\uparrow$ \\
OWL$^{\dagger}$        & one-shot pruning  & OPT-6.7B dense                & WikiText-2 perplexity $\downarrow$ \\
\bottomrule
\end{tabular}
\end{adjustbox}
\end{table}


The suite is ten frozen research repositories, listed in Table~\ref{tab:tasks} and chosen so that
between them they cover ten families of
training algorithm rather than ten instances of one: supervised fine-tuning, multi-turn agentic RL,
on-policy distillation, Bradley--Terry reward modeling, preference optimization, diffusion RL,
machine unlearning, discrete graph diffusion, weight averaging, and one-shot pruning. Each was
admitted on three properties the design depends on. It must ship a training algorithm its authors
actually run, not a tutorial or a toy; it must ship a frozen starting model, so that there is
something the agent is improving \emph{from}; and its metric must be recomputable, reproducibly, under
a half-day budget on a single B300, or the twelve-hour measurement in \S\ref{sec:protocol} could not
be run at all --- let alone once per cell.

Two of the ten do no training. Weight averaging combines a bank of checkpoints that are handed over
as data, and one-shot pruning removes weights from a released model in a single pass; neither has a
training horizon that could be extended or shortened. They are kept because the algorithmic question
is just as real in them --- which checkpoints to combine and how, which weights to remove and by what
criterion --- and because a suite that quietly dropped them would be a suite about training loops
rather than about algorithm design. The protocol treats them differently, and so does the baseline.

The ten metrics are incommensurable: an aesthetic score, a perplexity, a solve rate, a pass rate, an
unlearning balanced score. They cannot be averaged as they stand, and every per-task number in this
paper is reported in its own units against its own baseline.

\subsection{Protocol}
\label{sec:protocol}

Exploration is the same everywhere. For four hours the agent works inside the repository on one B300,
free to read it, edit it, launch training runs of its own, and consult the fast proxy metric without
limit. When the four hours are up it stops, and the code it leaves behind is the submission. Nothing
else crosses forward: no weights it trained, no cached state, no notes to itself --- only the source.

What happens next depends on the task. For the eight tasks that train, the submitted source is executed from
initialization until it terminates or twelve hours elapse, whichever comes first; the three most recent
checkpoints are then scored, and the cell takes the best of them under the task's direction. For the
two that do not train, the submitted code is simply executed once to produce its model --- an
averaged model, a pruned model --- which is scored directly; there the twelve hours are only a ceiling
on the verification stage, not a training budget being spent.

% The visibility boundary. This paragraph is deliberately explicit that the
% separation is one of access and timing, not of sample disjointness -- the honest
% version of the claim, and the one an auditor can check.
The boundary between what the agent may measure and what decides its score is the load-bearing
property of the whole design. During its four hours the agent may query the fast proxy as often as it
likes; the final metric is computed afterwards, from source it can no longer touch, by an evaluator
frozen before the first run. The separation is one of \emph{access and timing} rather than of sample
disjointness: on some tasks the cheap proxy is drawn from the same corpus the final evaluation uses,
because running the full evaluation as a proxy would cost more than the exploration budget allows. So
what the boundary guarantees is that no agent could score a candidate under the metric that decides
its result --- not that it never saw a row that metric would later read.

% Deleted on request: the two bookkeeping rules (a cell producing no runnable
% model is recorded as producing none rather than scored zero; missing values stay
% missing). Both still govern every rate in the paper, so if they are never stated
% they have to appear where the counts first appear, in \S\ref{sec:results}.

\subsection{Baselines}
\label{sec:baselines}

Calling a change an improvement requires something to compare it against, and the comparison this
paper reports is deliberately the strictest one available: the repository's own algorithm, given
exactly what the agent was given. For a task that trains, the baseline is the repository's committed
code executed under the identical procedure and the same twelve-hour budget, and measured by the same
fixed evaluator on the same asset. For the two tasks that do not train, it is the repository's recipe
executed as it stands, scored the same way. In both cases the only difference between the baseline and
a submission is the source code itself --- same hardware, same budget, same evaluator, same asset ---
which is what makes a win attributable to the change the agent made rather than to the resources it
was given.

% NOTE, to be resolved before submission: the agent's score is the best of its
% three most recent checkpoints (\S\ref{sec:protocol}). Whether the baseline is
% reduced the same way, or scored on a single final checkpoint, has to be stated
% here -- an earlier draft measured the two under different reduction rules and
% the pooled rate moved by around thirty points. Do not leave this implicit.
This is a harder bar than it may look, and a different one from what neighbouring benchmarks use. It
is not a published number, so it cannot have been tuned on a different evaluation than ours; it is not
an official instruction-tuned release, so beating it is not a matter of assembling more data than the
authors had; and it is not a human expert attempt, so it makes no claim about where human performance
lies. It is the answer to one question only: does the agent's code produce a better model than the
code that was already there, run under identical conditions?

\subsection{Scoring}
\label{sec:scoring}

Whether a submission beat the baseline is one bit, and one bit is too little for either use this
suite is meant to serve. Across tasks it makes a strong submission and a marginal one
indistinguishable, and it hides the difference between failing narrowly and failing completely. More
importantly, a benchmark of this shape is a natural environment for training agents by reinforcement
learning, and a binary outcome is a sparse reward: it gives no gradient between the many submissions
that do not beat the baseline, which on most tasks is where most submissions are. What is needed is a
dense score --- one that separates submissions everywhere along the range a metric can occupy, not
only at the point where it crosses a reference.

% The ladder, stated once as a formula. Everything task-specific lives in the
% triple (phi, worst, best); the score function itself never changes. An earlier
% design scaled by baseline +/- 2 sigma and saturated on the six tasks whose
% effects are far larger than their noise; this one uses no noise estimate and no
% property of the systems we happened to run.
Each task is therefore equipped with a \emph{progress coordinate} $\varphi$, a strictly increasing
function of quality that absorbs the metric's direction, together with three reference points in the
metric's own units: the uninformative model $x_{\perp}$, the baseline $x_{\mathrm{b}} = s(C)$, and the
optimum $x^{\ast}$. Writing $\phi_{\perp}, \phi_{\mathrm{b}}, \phi^{\ast}$ for their images under
$\varphi$, every task is scored by the same function,
\begin{equation}
  \label{eq:ladder}
  \sigma(x) \;=\;
  \begin{cases}
    \displaystyle 0.1\,\frac{\varphi(x) - \phi_{\perp}}{\phi_{\mathrm{b}} - \phi_{\perp}},
      & \varphi(x) \le \phi_{\mathrm{b}}, \\[1.1em]
    \displaystyle 0.1 + 0.9\,\frac{\varphi(x) - \phi_{\mathrm{b}}}{\phi^{\ast} - \phi_{\mathrm{b}}},
      & \varphi(x) > \phi_{\mathrm{b}},
  \end{cases}
\end{equation}
clipped to $[0,1]$, with $\sigma = 0$ for a submission that returned no model at all. The two branches
meet at $\sigma = 0.1$: matching the recipe the repository ships is the pivot of the scale, what lies
below it measures how far a submission fell short of that, and what lies above it measures how much
of the remaining distance to the optimum it closed.

Only the triple $(\varphi, x_{\perp}, x^{\ast})$ changes from task to task, and each element is fixed
by the metric rather than by the results. The optimum $x^{\ast}$ is the metric's best attainable
value: a rate of $1$, a preference score of $100$, a perplexity of $1$, a negative log-likelihood of
$0$. The uninformative point $x_{\perp}$ is what a model carrying no information about the task would
score --- $0$ for a rate, the chance level of $50$ for pairwise preference, the uniform predictor for
a likelihood --- which is why it is not always $0$ in the metric's own units. And $\varphi$ is the
identity wherever the metric is already a linear utility, which includes the rates, the aesthetic
head and the negative log-likelihood, and is $-\log$ for a perplexity, since a perplexity is the
exponential of a cross-entropy and only its logarithm lies on the same scale as the likelihood tasks.
Without that one transformation, a submission taking \textsc{owl} from $53.4$ to $16.2$ would read as
having closed $71\%$ of the distance to the optimum, where the correct figure is $30\%$: a perplexity
of $53.4$ is a cross-entropy of $\log 53.4 = 3.98$ nats above a perfect predictor and one of $16.2$ is
$2.79$ nats above it, so the submission removed $1.19$ of the $3.98$ nats that separated the
repository's own recipe from the optimum.

Every number in this paper that combines more than one task --- a system's average, the study-wide
mean --- is a mean of $\sigma$.
